\documentclass[11pt]{article}

\usepackage[margin=1in]{geometry}
\usepackage{amsmath,amssymb}
\usepackage{braket}
\usepackage{hyperref}

% Code formatting (portable: no external pygments needed)
\usepackage{xcolor}
\usepackage{listings}
\lstset{
  basicstyle=\ttfamily\small,
  keywordstyle=\color{blue!60!black},
  commentstyle=\color{green!40!black},
  stringstyle=\color{red!50!black},
  showstringspaces=false,
  breaklines=true,
  frame=single,
  columns=fullflexible
}

\title{Analytical (Parametric) Solutions of KL-Derived Linear Systems \\ \large Why Python/SymPy is Enough for 5--6 Qubits}
\author{}
\date{}

\begin{document}
\maketitle

\section{Problem description}

In many small quantum-code constructions (especially subset-sum / diagonal constructions), the Knill--Laflamme (KL) conditions for a chosen error set reduce to constraints on \emph{probabilities}
\[
X_j := |x_j|^2,\qquad Y_k := |y_k|^2,
\]
rather than directly on complex amplitudes $x_j,y_k$. This happens when all relevant operators are diagonal in the computational basis (e.g.\ single-qubit $Z_i$), so that expectation values depend only on squared magnitudes.

A typical setup looks like this (matching the user example).

\subsection*{Logical states supported on two classical sets}
We fix two sets of computational basis strings $C_0$ and $C_1$ and define
\[
\ket{0_L}=\sum_{x\in C_0} x_x \ket{x},\qquad
\ket{1_L}=\sum_{y\in C_1} y_y \ket{y}.
\]
Introduce probabilities
\[
X_x = |x_x|^2\ \ (x\in C_0), \qquad Y_y = |y_y|^2\ \ (y\in C_1),
\]
with normalization
\[
\sum_{x\in C_0} X_x = 1,\qquad \sum_{y\in C_1} Y_y = 1,
\qquad X_x\ge 0,\ Y_y\ge 0.
\]

\subsection*{KL constraints from diagonal errors}
For diagonal errors such as $Z_i$, the KL equal-expectation condition
\[
\langle 0_L|Z_i|0_L\rangle=\langle 1_L|Z_i|1_L\rangle
\]
becomes a linear equality in $X$'s and $Y$'s because $Z_i$ only contributes $\pm 1$ phases on basis states. Concretely, each expectation value is a signed sum of the relevant probabilities.

For later reference, if $\ket{\psi}=\sum_x c_x\ket{x}$ and $p_x:=|c_x|^2$, then
\begin{equation}
\label{eq:Zi-diagonal}
\bra{\psi}Z_i\ket{\psi}=\sum_x p_x\,(-1)^{x_i}.
\end{equation}


Therefore the full system has the structure
\[
A\mathbf{p} = \mathbf{b},\qquad \mathbf{p}\ge 0,
\]
together with two normalization constraints. Here $\mathbf{p}$ is the vector of all $X$'s and $Y$'s, and $A\mathbf{p}=\mathbf{b}$ collects the KL equalities.

\subsection*{What we want as an \emph{analytical general solution}}
A single feasible point is not enough. We want a \emph{parametric description}:
\[
\mathbf{p} = \mathbf{p}_0 + \sum_{r=1}^d t_r\, \mathbf{v}_r
\]
(with exact rational coefficients whenever possible), plus explicit inequalities describing the allowed region for the free parameters $t_r$ coming from $\mathbf{p}\ge 0$.

Geometrically, the solution set is a \emph{polytope}: an affine subspace cut by nonnegativity (simplex-type) inequalities.

\section{What Mathematica can do, and why we do not need it here}

Mathematica (Wolfram Language) is excellent at:
\begin{itemize}
  \item solving linear equalities symbolically (e.g.\ \texttt{Solve}),
  \item describing full solution sets of linear inequalities and equations (e.g.\ \texttt{Reduce}),
  \item optimizing linear objectives (e.g.\ \texttt{LinearOptimization} / \texttt{LinearProgramming}).
\end{itemize}

If you ask Mathematica for a general solution of a linear system with inequalities, \texttt{Reduce} can sometimes return a very clean description. This is especially convenient when you want projection/elimination without thinking about which variables to keep free.

\medskip
\noindent
\textbf{Why you do not need Mathematica here (for 5--6 qubits):}

\begin{itemize}
  \item The KL-derived systems you described are small and linear, and the coefficients are integers/rationals. Python's \texttt{sympy} does exact arithmetic over $\mathbb{Q}$ and gives analytical parametric solutions directly.
  \item The key task is not heavy optimization; it is \emph{elimination and parametrization}. For small systems we can do this explicitly: solve equalities, choose free variables, substitute, and read off inequalities.
  \item Mathematica's big advantage (strong automatic inequality elimination) becomes critical when the number of free parameters is large or you need a provably irredundant facet description. In your regime (5--6 qubits), it is usually unnecessary.
\end{itemize}

The main caveat is conceptual, not computational: this linear-probability method works when the KL constraints depend only on $|x|^2$ and $|y|^2$ (diagonal errors, disjoint support, etc.). If you enforce KL for errors that create off-diagonal overlaps (e.g.\ $X_i$-type errors), then constraints typically become \emph{quadratic} in amplitudes/phases and you leave the linear world. In that case you need polynomial methods (or relaxations), not just this pipeline.

\section{What we do in Python (the method)}

The Python approach is a direct algebraic pipeline:

\subsection*{Step A: Build the linear equalities}
Write all KL equalities and the two normalization equations as linear equations in variables
\[
X_0,\dots,X_{|C_0|-1},\qquad Y_0,\dots,Y_{|C_1|-1}.
\]

\subsection*{Step B: Solve equalities exactly (symbolic, rational)}
Use SymPy's linear solving (via \texttt{solve} or matrix methods) to express a set of dependent variables in terms of chosen free variables. This yields an affine parametrization.

\subsection*{Step C: Impose nonnegativity by substitution}
Nonnegativity constraints $X_j\ge 0$ and $Y_k\ge 0$ become linear inequalities in the free parameters after substitution. Collect and simplify them.

\subsection*{Step D: Present as a general analytical solution}
The output is:
\begin{itemize}
  \item explicit formulas for all $X$'s and $Y$'s in terms of a small number of free parameters,
  \item explicit linear constraints on those free parameters (e.g.\ $t_1\ge0$, $t_2\ge0$, $t_1+t_2\le 2/7$),
  \item optionally, auto-generated \LaTeX\ for both formulas and inequalities.
\end{itemize}

\section{Worked example (matching the provided system)}

In your concrete example there are
\[
(X_0,X_1,X_2)\quad \text{and}\quad (Y_0,\dots,Y_5)
\]
with KL constraints from single-qubit $Z_i$ giving five linear equations, plus two normalization equations.

Solving the equalities yields the affine description
\[
X_0 = \frac{3}{7},\quad X_1=X_2=\frac{2}{7},
\]
\[
Y_0 = \frac{1}{7}+Y_5,\quad Y_1=\frac{1}{7}+Y_4,\quad
Y_2=\frac{3}{7}-Y_4-Y_5,\quad Y_3=\frac{2}{7}-Y_4-Y_5,
\]
and the inequalities $Y_k\ge 0$ reduce to
\[
Y_4\ge 0,\quad Y_5\ge 0,\quad Y_4+Y_5\le \frac{2}{7}.
\]
This is exactly the kind of ``general solution'' you want, and it is obtained by linear algebra plus inequality substitution.


\section{Five-qubit examples: Families A--E (from \texttt{lambdav2.tex})}

The file \texttt{lambdav2.tex} contains five explicit five-qubit constructions (Families A--E) built from subset-sum supports.
For the purpose of this note, the subset-sum origin is not essential: once the two supports $C_0$ and $C_{S_1}$ are fixed, the ``shared $Z$-marginal'' KL constraints become a small linear feasibility problem in the probabilities $\{X_x\}_{x\in C_0}$ and $\{Y_y\}_{y\in C_{S_1}}$.

In each family we impose:
\[
\sum_{x\in C_0} X_x = 1,\qquad \sum_{y\in C_{S_1}} Y_y = 1,\qquad X_x\ge 0,\ Y_y\ge 0,
\]
and for $i=1,\dots,5$,
\[
\sum_{x\in C_0} (-1)^{x_i} X_x \;=\; \sum_{y\in C_{S_1}} (-1)^{y_i} Y_y,
\]
which is exactly the condition $\langle 0_L|Z_i|0_L\rangle=\langle 1_L|Z_i|1_L\rangle$ written at the level of squared amplitudes.

For Families A--E below, the affine solution space always has dimension $1$, so the feasible set is a line segment; this explains why \texttt{lambdav2.tex} presents each family using a single parameter $\lambda\in[0,1]$.
For comparison, the worked example in Section~4 has \emph{two} free parameters.


\begin{table}[h]
\centering
\begin{tabular}{c|c|c|c|c|c|c|c}
Family & $m$ & $\mathbf w$ & $(S_0,S_1)$ & $|C_0|$ & $|C_{S_1}|$ & \#free & free param range \\
\hline
A & 6 & $(4,5,1,2,2)$ & $(0,3)$ & 3 & 3 & 1 & $Y_{00110}\in[0,\frac{1}{2}]$ \\
B & 6 & $(1,1,1,5,3)$ & $(0,2)$ & 4 & 4 & 1 & $Y_{11000}\in[0,\frac{1}{3}]$ \\
C & 6 & $(3,3,1,1,1)$ & $(0,4)$ & 4 & 4 & 1 & $Y_{10100}\in[0,\frac{1}{3}]$ \\
D & 6 & $(3,1,3,5,5)$ & $(0,4)$ & 4 & 4 & 1 & $Y_{10111}\in[0,\frac{1}{3}]$ \\
E & 8 & $(1,6,2,5,5)$ & $(0,4)$ & 3 & 3 & 1 & $Y_{11001}\in[0,\frac{1}{2}]$ \\
\end{tabular}
\caption{Summary of the five-qubit families A--E: support sizes, linear-system nullity, and the resulting one-parameter feasibility interval for a convenient free probability.}
\end{table}

\subsection{Family A: $(m,\mathbf w,S_0,S_1)=(6,(4,5,1,2,2),0,3)$}

\paragraph{Supports.}
\[
C_0 = \{10010,\ 11101,\ 01100\},\qquad C_3 = \{11000,\ 10111,\ 00110\}.
\]

\paragraph{Variables.}
For each $x\in C_0$ let $X_x:=|x_x|^2$, and for each $y\in C_3$ let $Y_y:=|y_y|^2$.
Normalization and nonnegativity give $\sum_{x\in C_0}X_x=1$, $\sum_{y\in C_3}Y_y=1$, and all $X_x,Y_y\ge 0$.

\paragraph{KL constraints from $Z_i$.}
For $i=1,\dots,5$ we impose the shared $Z$-marginals
\[
\langle 0_L|Z_i|0_L\rangle=\langle 1_L|Z_i|1_L\rangle,
\]
which become linear equations in $\{X_x\}$ and $\{Y_y\}$ because $Z_i\ket{x}=(-1)^{x_i}\ket{x}$.
The resulting linear system has 6 variables and rank 5, hence a 1-parameter affine solution set. After intersecting with nonnegativity, the feasible set is a convex polytope (here a line segment).

\paragraph{General analytical solution.}
\[
\begin{aligned}
X_{10010} = Y_{11000} &= \frac{1}{2}\\
X_{11101} = Y_{10111} &= \frac{1}{2} - t\\
X_{01100} = Y_{00110} &= t\\
\text{with }&\quad 0 \leq t \wedge t \leq \frac{1}{2}.
\end{aligned}
\]
\paragraph{One convenient choice of logical states.}
Choosing real nonnegative amplitudes (phases can be added without affecting these $Z_i$ constraints), one may take
\[
\ket{0_L(t)} = \sqrt{\frac{1}{2}}\ket{10010} + \sqrt{\frac{1}{2} - t}\ket{11101} + \sqrt{t}\ket{01100},\qquad \ket{1_L(t)} = \sqrt{\frac{1}{2}}\ket{11000} + \sqrt{\frac{1}{2} - t}\ket{10111} + \sqrt{t}\ket{00110}.
\]
\paragraph{Relation to the $\lambda$ parametrization in \texttt{lambdav2.tex}.}
In \texttt{lambdav2.tex} the family is written with $\lambda\in[0,1]$ so that the third probability equals $\lambda/2$. This corresponds to the reparametrization $t=\lambda/2$.


\subsection{Family B: $(m,\mathbf w,S_0,S_1)=(6,(1,1,1,5,3),0,2)$}

\paragraph{Supports.}
\[
C_0 = \{00000,\ 00110,\ 10010,\ 11101\},\qquad C_2 = \{00011,\ 01100,\ 10100,\ 11000\}.
\]

\paragraph{Variables.}
For each $x\in C_0$ let $X_x:=|x_x|^2$, and for each $y\in C_2$ let $Y_y:=|y_y|^2$.
Normalization and nonnegativity give $\sum_{x\in C_0}X_x=1$, $\sum_{y\in C_2}Y_y=1$, and all $X_x,Y_y\ge 0$.

\paragraph{KL constraints from $Z_i$.}
For $i=1,\dots,5$ we impose the shared $Z$-marginals
\[
\langle 0_L|Z_i|0_L\rangle=\langle 1_L|Z_i|1_L\rangle,
\]
which become linear equations in $\{X_x\}$ and $\{Y_y\}$ because $Z_i\ket{x}=(-1)^{x_i}\ket{x}$.
The resulting linear system has 8 variables and rank 7, hence a 1-parameter affine solution set. After intersecting with nonnegativity, the feasible set is a convex polytope (here a line segment).

\paragraph{General analytical solution.}
\[
\begin{aligned}
X_{00000} = X_{11101} = Y_{00011} = Y_{10100} &= \frac{1}{3}\\
X_{00110} = Y_{01100} &= \frac{1}{3} - t\\
X_{10010} = Y_{11000} &= t\\
\text{with }&\quad 0 \leq t \wedge t \leq \frac{1}{3}.
\end{aligned}
\]
\paragraph{One convenient choice of logical states.}
Choosing real nonnegative amplitudes (phases can be added without affecting these $Z_i$ constraints), one may take
\[
\ket{0_L(t)} = \sqrt{\frac{1}{3}}\ket{00000} + \sqrt{\frac{1}{3} - t}\ket{00110} + \sqrt{t}\ket{10010} + \sqrt{\frac{1}{3}}\ket{11101},\qquad \ket{1_L(t)} = \sqrt{\frac{1}{3}}\ket{00011} + \sqrt{\frac{1}{3} - t}\ket{01100} + \sqrt{\frac{1}{3}}\ket{10100} + \sqrt{t}\ket{11000}.
\]
\paragraph{Relation to the $\lambda$ parametrization in \texttt{lambdav2.tex}.}
In \texttt{lambdav2.tex} the family is written with $\lambda\in[0,1]$ so that one of the varying probabilities equals $\lambda/3$ and the complementary one equals $(1-\lambda)/3$. With the present choice of free parameter $t$ (the $(1-\lambda)/3$ branch), one has $t=(1-\lambda)/3$.


\subsection{Family C: $(m,\mathbf w,S_0,S_1)=(6,(3,3,1,1,1),0,4)$}

\paragraph{Supports.}
\[
C_0 = \{00000,\ 01111,\ 10111,\ 11000\},\qquad C_4 = \{01001,\ 01100,\ 10010,\ 10100\}.
\]

\paragraph{Variables.}
For each $x\in C_0$ let $X_x:=|x_x|^2$, and for each $y\in C_4$ let $Y_y:=|y_y|^2$.
Normalization and nonnegativity give $\sum_{x\in C_0}X_x=1$, $\sum_{y\in C_4}Y_y=1$, and all $X_x,Y_y\ge 0$.

\paragraph{KL constraints from $Z_i$.}
For $i=1,\dots,5$ we impose the shared $Z$-marginals
\[
\langle 0_L|Z_i|0_L\rangle=\langle 1_L|Z_i|1_L\rangle,
\]
which become linear equations in $\{X_x\}$ and $\{Y_y\}$ because $Z_i\ket{x}=(-1)^{x_i}\ket{x}$.
The resulting linear system has 8 variables and rank 7, hence a 1-parameter affine solution set. After intersecting with nonnegativity, the feasible set is a convex polytope (here a line segment).

\paragraph{General analytical solution.}
\[
\begin{aligned}
X_{00000} = X_{11000} = Y_{01001} = Y_{10010} &= \frac{1}{3}\\
X_{01111} = Y_{01100} &= \frac{1}{3} - t\\
X_{10111} = Y_{10100} &= t\\
\text{with }&\quad 0 \leq t \wedge t \leq \frac{1}{3}.
\end{aligned}
\]
\paragraph{One convenient choice of logical states.}
Choosing real nonnegative amplitudes (phases can be added without affecting these $Z_i$ constraints), one may take
\[
\ket{0_L(t)} = \sqrt{\frac{1}{3}}\ket{00000} + \sqrt{\frac{1}{3} - t}\ket{01111} + \sqrt{t}\ket{10111} + \sqrt{\frac{1}{3}}\ket{11000},\qquad \ket{1_L(t)} = \sqrt{\frac{1}{3}}\ket{01001} + \sqrt{\frac{1}{3} - t}\ket{01100} + \sqrt{\frac{1}{3}}\ket{10010} + \sqrt{t}\ket{10100}.
\]
\paragraph{Relation to the $\lambda$ parametrization in \texttt{lambdav2.tex}.}
In \texttt{lambdav2.tex} the family is written with $\lambda\in[0,1]$ so that one of the varying probabilities equals $\lambda/3$ and the complementary one equals $(1-\lambda)/3$. With the present choice of free parameter $t$ (the $(1-\lambda)/3$ branch), one has $t=(1-\lambda)/3$.


\subsection{Family D: $(m,\mathbf w,S_0,S_1)=(6,(3,1,3,5,5),0,4)$}

\paragraph{Supports.}
\[
C_0 = \{00000,\ 01001,\ 11101,\ 11110\},\qquad C_4 = \{00011,\ 01100,\ 10111,\ 11000\}.
\]

\paragraph{Variables.}
For each $x\in C_0$ let $X_x:=|x_x|^2$, and for each $y\in C_4$ let $Y_y:=|y_y|^2$.
Normalization and nonnegativity give $\sum_{x\in C_0}X_x=1$, $\sum_{y\in C_4}Y_y=1$, and all $X_x,Y_y\ge 0$.

\paragraph{KL constraints from $Z_i$.}
For $i=1,\dots,5$ we impose the shared $Z$-marginals
\[
\langle 0_L|Z_i|0_L\rangle=\langle 1_L|Z_i|1_L\rangle,
\]
which become linear equations in $\{X_x\}$ and $\{Y_y\}$ because $Z_i\ket{x}=(-1)^{x_i}\ket{x}$.
The resulting linear system has 8 variables and rank 7, hence a 1-parameter affine solution set. After intersecting with nonnegativity, the feasible set is a convex polytope (here a line segment).

\paragraph{General analytical solution.}
\[
\begin{aligned}
X_{00000} = X_{11110} = Y_{01100} = Y_{11000} &= \frac{1}{3}\\
X_{01001} = Y_{00011} &= \frac{1}{3} - t\\
X_{11101} = Y_{10111} &= t\\
\text{with }&\quad 0 \leq t \wedge t \leq \frac{1}{3}.
\end{aligned}
\]
\paragraph{One convenient choice of logical states.}
Choosing real nonnegative amplitudes (phases can be added without affecting these $Z_i$ constraints), one may take
\[
\ket{0_L(t)} = \sqrt{\frac{1}{3}}\ket{00000} + \sqrt{\frac{1}{3} - t}\ket{01001} + \sqrt{t}\ket{11101} + \sqrt{\frac{1}{3}}\ket{11110},\qquad \ket{1_L(t)} = \sqrt{\frac{1}{3} - t}\ket{00011} + \sqrt{\frac{1}{3}}\ket{01100} + \sqrt{t}\ket{10111} + \sqrt{\frac{1}{3}}\ket{11000}.
\]
\paragraph{Relation to the $\lambda$ parametrization in \texttt{lambdav2.tex}.}
In \texttt{lambdav2.tex} the family is written with $\lambda\in[0,1]$ so that one of the varying probabilities equals $\lambda/3$ and the complementary one equals $(1-\lambda)/3$. With the present choice of free parameter $t$ (the $(1-\lambda)/3$ branch), one has $t=(1-\lambda)/3$.


\subsection{Family E: $(m,\mathbf w,S_0,S_1)=(8,(1,6,2,5,5),0,4)$}

\paragraph{Supports.}
\[
C_0 = \{01011,\ 10110,\ 10101\},\qquad C_4 = \{00111,\ 11010,\ 11001\}.
\]

\paragraph{Variables.}
For each $x\in C_0$ let $X_x:=|x_x|^2$, and for each $y\in C_4$ let $Y_y:=|y_y|^2$.
Normalization and nonnegativity give $\sum_{x\in C_0}X_x=1$, $\sum_{y\in C_4}Y_y=1$, and all $X_x,Y_y\ge 0$.

\paragraph{KL constraints from $Z_i$.}
For $i=1,\dots,5$ we impose the shared $Z$-marginals
\[
\langle 0_L|Z_i|0_L\rangle=\langle 1_L|Z_i|1_L\rangle,
\]
which become linear equations in $\{X_x\}$ and $\{Y_y\}$ because $Z_i\ket{x}=(-1)^{x_i}\ket{x}$.
The resulting linear system has 6 variables and rank 5, hence a 1-parameter affine solution set. After intersecting with nonnegativity, the feasible set is a convex polytope (here a line segment).

\paragraph{General analytical solution.}
\[
\begin{aligned}
X_{01011} = Y_{00111} &= \frac{1}{2}\\
X_{10110} = Y_{11010} &= \frac{1}{2} - t\\
X_{10101} = Y_{11001} &= t\\
\text{with }&\quad 0 \leq t \wedge t \leq \frac{1}{2}.
\end{aligned}
\]
\paragraph{One convenient choice of logical states.}
Choosing real nonnegative amplitudes (phases can be added without affecting these $Z_i$ constraints), one may take
\[
\ket{0_L(t)} = \sqrt{\frac{1}{2}}\ket{01011} + \sqrt{\frac{1}{2} - t}\ket{10110} + \sqrt{t}\ket{10101},\qquad \ket{1_L(t)} = \sqrt{\frac{1}{2}}\ket{00111} + \sqrt{\frac{1}{2} - t}\ket{11010} + \sqrt{t}\ket{11001}.
\]
\paragraph{Relation to the $\lambda$ parametrization in \texttt{lambdav2.tex}.}
In \texttt{lambdav2.tex} the family is written with $\lambda\in[0,1]$ so that the third probability equals $\lambda/2$. This corresponds to the reparametrization $t=\lambda/2$.


\section{Python code (SymPy) to obtain analytical general solutions}

This section contains self-contained Python code that:
\begin{itemize}
  \item builds the KL ``shared $Z$-marginal'' equations from two supports $C_0$ and $C_{S_1}$,
  \item solves the linear equalities exactly (over $\mathbb{Q}$) using \texttt{sympy},
  \item identifies the true free probability parameters (the affine nullspace),
  \item derives the allowed parameter region coming from nonnegativity.
\end{itemize}

\noindent
\textbf{Remark.} This code solves the \emph{linear-probability} KL constraints (diagonal errors).
If you add KL constraints for non-diagonal errors (e.g.\ $X_i$), the constraints become quadratic in amplitudes/phases and this LP-style pipeline no longer applies.

\begin{lstlisting}[language=Python]
import sympy as sp
from sympy.solvers.inequalities import reduce_inequalities

def z_sign(bitstring, i):
    # eigenvalue of Z_i on |bitstring> ( +1 for 0, -1 for 1 )
    return 1 if bitstring[i] == "0" else -1

def build_Z_marginal_system(C0, C1):
    # Build linear equalities:
    #   <0_L|Z_i|0_L> = <1_L|Z_i|1_L> for i=1..n,
    # plus separate normalizations:
    #   sum_{x in C0} X_x = 1,  sum_{y in C1} Y_y = 1.
    n = len(C0[0])
    X = sp.symbols(f"X0:{len(C0)}", real=True)
    Y = sp.symbols(f"Y0:{len(C1)}", real=True)
    vars_all = list(X) + list(Y)

    eqs = []
    for i in range(n):
        lhs = sum(X[j] * z_sign(C0[j], i) for j in range(len(C0)))
        rhs = sum(Y[k] * z_sign(C1[k], i) for k in range(len(C1)))
        eqs.append(sp.Eq(lhs, rhs))

    eqs.append(sp.Eq(sum(X), 1))
    eqs.append(sp.Eq(sum(Y), 1))

    return X, Y, vars_all, eqs

def solve_general_probabilities(C0, C1):
    # Return an exact affine solution set for probabilities, plus nonnegativity constraints.
    # Output:
    #   mapping   : dict var -> affine expression in free variables
    #   free_vars : list of free variables (a subset of vars_all) left by linsolve
    #   region    : if len(free_vars)==1, a simplified interval constraint; otherwise a list of inequalities

    X, Y, vars_all, eqs = build_Z_marginal_system(C0, C1)

    sol_set = sp.linsolve(eqs, vars_all)
    if not sol_set:
        raise ValueError("No solutions: the linear constraints are inconsistent.")

    sol_tuple = list(sol_set)[0]
    mapping = {v: sp.simplify(expr) for v, expr in zip(vars_all, sol_tuple)}

    # SymPy represents free parameters either by leaving some variables unchanged
    # (mapping[v] == v) or by introducing internal parameters.
    free_vars = [v for v in vars_all if sp.simplify(mapping[v] - v) == 0]

    # Nonnegativity constraints: all probabilities >= 0
    ineqs = [sp.Ge(mapping[v], 0) for v in vars_all]
    ineqs = [sp.simplify(i) for i in ineqs if i != True]

    if len(free_vars) == 0:
        # Unique solution: either feasible or infeasible (check inequalities).
        region = True if all(sp.simplify(i) != False for i in ineqs) else False
        return mapping, free_vars, region

    if len(free_vars) == 1:
        # In 1D, reduce to a clean interval like 0 <= t <= 1/3.
        t = free_vars[0]
        region = sp.simplify(reduce_inequalities(ineqs, t))
        return mapping, free_vars, region

    # For >=2 parameters, return the inequality list (still exact and linear).
    return mapping, free_vars, ineqs

def pretty_print_solution(C0, C1, mapping, free_vars, region):
    # Print the solution in terms of codewords rather than X0,X1,... and Y0,Y1,...
    X_syms = [s for s in mapping.keys() if str(s).startswith("X")]
    Y_syms = [s for s in mapping.keys() if str(s).startswith("Y")]
    X_syms = sorted(X_syms, key=lambda s: int(str(s)[1:]))  # X0,X1,...
    Y_syms = sorted(Y_syms, key=lambda s: int(str(s)[1:]))  # Y0,Y1,...

    print("Free variables:", free_vars)
    print("Region:", region)
    print("Probabilities on C0:")
    for j, x in enumerate(C0):
        print(f"  X_{{{x}}} = {mapping[X_syms[j]]}")
    print("Probabilities on C1:")
    for k, y in enumerate(C1):
        print(f"  Y_{{{y}}} = {mapping[Y_syms[k]]}")

# --------------------------
# Examples: Families A--E (supports taken from lambdav2.tex)
# --------------------------
families = {
    "A": (["10010","11101","01100"], ["11000","10111","00110"]),
    "B": (["00000","00110","10010","11101"], ["00011","01100","10100","11000"]),
    "C": (["00000","01111","10111","11000"], ["01001","01100","10010","10100"]),
    "D": (["00000","01001","11101","11110"], ["00011","01100","10111","11000"]),
    "E": (["01011","10110","10101"], ["00111","11010","11001"]),
}

for name, (C0, C1) in families.items():
    print("\\n=== Family", name, "===")
    mapping, free_vars, region = solve_general_probabilities(C0, C1)
    pretty_print_solution(C0, C1, mapping, free_vars, region)
\end{lstlisting}

\section{How to use this at scale}

For many instances (many small codes), the scalable pattern is:

\begin{itemize}
  \item Generate the KL linear equalities automatically from your classical sets $C_0,C_1$ and your error list (e.g.\ $Z_i$).
  \item Choose free parameters in a consistent, human-readable way (for example, keep the last few $Y$'s free).
  \item Call \texttt{affine\_polytope\_solution} to get exact formulas and constraints.
  \item Emit \LaTeX\ by \texttt{latex\_report} and paste into your writeup, or auto-write a \texttt{.tex} snippet per code family.
\end{itemize}

In the common situation where the number of free parameters is only 1--3, the resulting inequalities are typically already as clean as the Mathematica \texttt{Reduce} output, without needing any heavier machinery.

\section{Summary}

When KL conditions reduce to linear constraints on probabilities, the ``general analytical solution'' is a polytope description: an affine parametrization plus linear inequalities. For 5--6 qubits, Python + SymPy is enough to derive these solutions exactly (over rationals), produce compact parametric families, and export \LaTeX. Mathematica is convenient but not necessary in this regime.



%%%%%%%%%%%%%%%%%%%%%%%%%%%%%%%%%%%%%%%%%%%%%%%%%%%%%%%%%%%%%%%%%%%%%%%%%%%%%%%%
\appendix
\section{Appendix: How to solve these KL-derived LPs analytically (and why $\mathrm{conv}(V_i)$ matters)}
\label{sec:appendix-procedure}

This appendix expands on the workflow discussed in the chat: in the distance-$2$ SSLP regime, the KL constraints for single-qubit $Z_i$ errors become a \emph{linear feasibility problem} in the probability variables. The ``analytical general solution'' is obtained by (i) exact linear algebra over $\mathbb{Q}$ (Gaussian elimination / RREF) and then (ii) imposing nonnegativity, which cuts out a convex polytope in the free-parameter space.

\subsection{Step-by-step analytical procedure}

Fix two classical supports $C_0$ and $C_1$ (coming from subset-sum classes, or any other construction). Write
\[
\ket{0_L}=\sum_{x\in C_0} a_x\ket{x},\qquad
\ket{1_L}=\sum_{y\in C_1} b_y\ket{y},
\]
with probabilities
\[
X_x:=|a_x|^2\ (x\in C_0),\qquad Y_y:=|b_y|^2\ (y\in C_1),
\]
subject to
\[
\sum_{x\in C_0}X_x=1,\qquad \sum_{y\in C_1}Y_y=1,\qquad X_x\ge 0,\ Y_y\ge 0.
\]

\paragraph{Step 1 (encode $Z_i$ by sign patterns).}
For a bitstring $x\in\{0,1\}^n$, define its \emph{$Z$-sign vector}
\[
 s(x)\in\{\pm 1\}^n,\qquad s_i(x):=(-1)^{x_i}.
\]
This records the eigenvalues of the single-site operators $Z_i$ on $\ket{x}$.

\paragraph{Step 2 (write the shared $Z$-marginal equalities).}
Using $Z_i\ket{x}=(-1)^{x_i}\ket{x}$ and the diagonal formula \eqref{eq:Zi-diagonal}, the KL equalities
\[
\bra{0_L}Z_i\ket{0_L}=\bra{1_L}Z_i\ket{1_L}\qquad (i=1,\dots,n)
\]
become
\[
\sum_{x\in C_0} s_i(x)\,X_x\;=\;\sum_{y\in C_1} s_i(y)\,Y_y\qquad (i=1,\dots,n).
\]
Together with the two normalization equations, these are linear equalities over $\mathbb{Q}$.

\paragraph{Step 3 (matrix form).}
Order the supports as $C_0=\{x^{(1)},\dots,x^{(N_0)}\}$ and $C_1=\{y^{(1)},\dots,y^{(N_1)}\}$. Form the integer matrices
\[
M_0:=\big[s(x^{(1)})\ \cdots\ s(x^{(N_0)})\big]\in\mathbb{Z}^{n\times N_0},
\qquad
M_1:=\big[s(y^{(1)})\ \cdots\ s(y^{(N_1)})\big]\in\mathbb{Z}^{n\times N_1}.
\]
Let $X\in\mathbb{R}^{N_0}$ and $Y\in\mathbb{R}^{N_1}$ be the probability vectors. Then
\[
\big(\bra{0_L}Z_1\ket{0_L},\dots,\bra{0_L}Z_n\ket{0_L}\big)=M_0X,
\qquad
\big(\bra{1_L}Z_1\ket{1_L},\dots,\bra{1_L}Z_n\ket{1_L}\big)=M_1Y.
\]
The shared-marginal constraints are
\begin{equation}
\label{eq:matrix-ZKL}
M_0X=M_1Y,
\end{equation}
plus
\begin{equation}
\label{eq:matrix-norm}
\mathbf{1}^TX=1,\qquad \mathbf{1}^TY=1.
\end{equation}
Equivalently, stack $p:=(X,Y)\in\mathbb{R}^{N_0+N_1}$ and write
\[
A p=b
\]
with $A\in\mathbb{Z}^{(n+2)\times(N_0+N_1)}$.

\paragraph{Step 4 (Gaussian elimination / RREF over $\mathbb{Q}$).}
Compute the row-reduced echelon form of the augmented matrix $[A\mid b]$ over exact rationals. This yields an \emph{affine solution space}
\[
 p=p_{\mathrm{part}}+\sum_{r=1}^d t_r v_r,
\]
where $d$ is the number of free parameters (the nullity of $A$).

\paragraph{Step 5 (impose nonnegativity: a polytope in parameter space).}
Substitute the affine parametrization into the inequalities $p\ge 0$. Each coordinate inequality becomes linear in $(t_1,\dots,t_d)$:
\[
\alpha_0+\sum_{r=1}^d \alpha_r t_r\ge 0.
\]
Hence the feasible parameter set is a convex polytope
\[
\mathcal{P}:=\{t\in\mathbb{R}^d: Bt\le c\}.
\]
This is the step where ``is the convex region empty?'' is decided: the system is feasible iff $\mathcal{P}\neq\varnothing$.

\paragraph{Step 6 (simplify inequalities).}
In small dimensions ($d\le 2$ or $3$), one usually simplifies by hand:
\begin{itemize}
\item In $d=1$, all constraints reduce to an interval $t\in[t_{\min},t_{\max}]$.
\item In $d=2$, the region is a polygon; redundant constraints can often be dropped by inspection.
\item Systematically, one can use Fourier--Motzkin elimination or a vertex computation, but for $5$--$6$ qubits it is typically overkill.
\end{itemize}

\paragraph{Step 7 (lift back to amplitudes).}
Once probabilities are fixed, one convenient normalized choice is to take real nonnegative amplitudes
\[
\ket{0_L(t)}=\sum_{x\in C_0} \sqrt{X_x(t)}\,\ket{x},\qquad
\ket{1_L(t)}=\sum_{y\in C_1} \sqrt{Y_y(t)}\,\ket{y}.
\]
Phases $e^{i\phi_x},e^{i\theta_y}$ may be inserted if desired; they do not affect the diagonal $Z_i$ expectations.

\subsection{The convex-hull viewpoint and $\mathrm{conv}(V_i)$}
\label{subsec:convVi}

The same linear feasibility problem has a very useful geometric interpretation.

\paragraph{Definition of $V_0,V_1$ and their convex hulls.}
From the supports define the finite vertex sets
\[
V_0:=\{s(x):x\in C_0\}\subseteq\{\pm 1\}^n,\qquad
V_1:=\{s(y):y\in C_1\}\subseteq\{\pm 1\}^n.
\]
Their convex hulls are
\[
\mathrm{conv}(V_0)=\Big\{\sum_{x\in C_0} X_x\,s(x): X_x\ge 0,\ \sum_{x\in C_0}X_x=1\Big\},
\]
\[
\mathrm{conv}(V_1)=\Big\{\sum_{y\in C_1} Y_y\,s(y): Y_y\ge 0,\ \sum_{y\in C_1}Y_y=1\Big\}.
\]
Thus $\mathrm{conv}(V_i)$ is exactly the set of all achievable single-site $Z$-expectation vectors from probability distributions supported on $C_i$.

\paragraph{Why it is equivalent to the $Z$-KL equalities.}
Define
\[
 t_0:=\big(\bra{0_L}Z_1\ket{0_L},\dots,\bra{0_L}Z_n\ket{0_L}\big)=M_0X\in\mathrm{conv}(V_0),
\]
\[
 t_1:=\big(\bra{1_L}Z_1\ket{1_L},\dots,\bra{1_L}Z_n\ket{1_L}\big)=M_1Y\in\mathrm{conv}(V_1).
\]
The shared-marginal constraints $t_0=t_1$ are therefore equivalent to the existence of a common vector
\[
 t\in \mathrm{conv}(V_0)\cap \mathrm{conv}(V_1).
\]

\paragraph{When do the shared $Z$-marginals vary with LP parameters?}
Two different phenomena can occur:
\begin{itemize}
\item The affine feasibility set in $(X,Y)$ can have positive dimension (free parameters), but the induced marginal vector $t$ can still be fixed.
\item The marginal vector $t$ varies with parameters exactly when $\mathrm{conv}(V_0)\cap\mathrm{conv}(V_1)$ has positive dimension (e.g. contains a line segment).
\end{itemize}
The examples below illustrate both cases.

\subsection{Example walk-through 1: Family A (one free parameter and \emph{varying} $Z$-marginals)}
\label{subsec:walkthrough-familyA}

Recall Family~A in Section~5:
\[
C_0=\{10010,11101,01100\},\qquad C_3=\{11000,10111,00110\}.
\]
Introduce variables
\[
X_1:=X_{10010},\ X_2:=X_{11101},\ X_3:=X_{01100},\qquad
Y_1:=Y_{11000},\ Y_2:=Y_{10111},\ Y_3:=Y_{00110}.
\]

\paragraph{Step A1: write the $Z$-sign vectors.}
Using $s_i(x)=(-1)^{x_i}$, we have
\[
\begin{aligned}
 s(10010)&=(-1,+1,+1,-1,+1),\\
 s(11101)&=(-1,-1,-1,+1,-1),\\
 s(01100)&=(+1,-1,-1,+1,+1),
\end{aligned}
\qquad
\begin{aligned}
 s(11000)&=(-1,-1,+1,+1,+1),\\
 s(10111)&=(-1,+1,-1,-1,-1),\\
 s(00110)&=(+1,+1,-1,-1,+1).
\end{aligned}
\]

\paragraph{Step A2: write the equalities and do elimination.}
The shared-marginal equalities are
\[
\sum_{j=1}^3 s_i(x^{(j)})X_j=\sum_{k=1}^3 s_i(y^{(k)})Y_k\qquad (i=1,\dots,5),
\]
together with
\(
X_1+X_2+X_3=1
\) and
\(
Y_1+Y_2+Y_3=1.
\)
A convenient elimination route is:
\begin{itemize}
\item From the $i=3$ and $i=4$ equations one obtains $X_1=X_2+X_3$.
\item Using $X_1+X_2+X_3=1$ then gives $X_1=\tfrac12$ and $X_2+X_3=\tfrac12$.
\item Similarly, combining $i=2$ and $i=3$ yields $Y_1=Y_2+Y_3$, and then $Y_1=\tfrac12$.
\item Finally, the $i=5$ equation forces $X_3=Y_3$.
\end{itemize}
Thus letting the free parameter be
\(
 t:=X_3=Y_3
\), we obtain the affine solution
\[
X_1=Y_1=\frac12,\qquad X_2=Y_2=\frac12-t,\qquad X_3=Y_3=t.
\]

\paragraph{Step A3: impose nonnegativity.}
The constraints $X_2\ge 0$ and $X_3\ge 0$ give
\[
0\le t\le \frac12.
\]

\paragraph{Step A4: compute the shared $Z$-marginals and their range.}
Using \eqref{eq:Zi-diagonal} (or directly $t=\sum p\,s$), one finds
\[
\bigl(\langle Z_1\rangle,\dots,\langle Z_5\rangle\bigr)=\bigl(2t-1,\ 0,\ 0,\ 0,\ 2t\bigr),
\]
so the common marginal vector varies with $t$.
In particular,
\[
\Sigma_Z(t):=\sum_{i=1}^5\langle Z_i\rangle^2=(2t-1)^2+(2t)^2=8t^2-4t+1
\]
ranges over
\[
\Sigma_Z(t)\in\left[\frac12,\ 1\right]\quad\text{for }t\in\left[0,\frac12\right],
\]
with minimum at $t=\tfrac14$.

\paragraph{Step A5: interpret via $\mathrm{conv}(V_0)\cap\mathrm{conv}(V_1)$.}
Here the intersection $\mathrm{conv}(V_0)\cap\mathrm{conv}(V_1)$ contains a line segment (parameterized by $t$), which is why the shared $Z$-marginals move continuously.

\subsection{Example walk-through 2: the worked $m=7$ example (two free parameters but \emph{fixed} $Z$-marginals)}
\label{subsec:walkthrough-worked}

Section~4 stated a two-parameter general solution with inequalities
\[
Y_4\ge 0,\quad Y_5\ge 0,\quad Y_4+Y_5\le \frac{2}{7}.
\]
For completeness we record the underlying supports (as in the worked $((5,2,2))$ example discussed earlier):
\[
C_0=\{00000,01111,10111\},
\qquad
C_1=\{00011,00101,00110,11001,11010,11100\}.
\]
The affine solution is
\[
X_{00000}=\frac{3}{7},\quad X_{01111}=X_{10111}=\frac{2}{7},
\]
\[
\begin{aligned}
Y_{00011}&=\frac{1}{7}+Y_{11100},\qquad
Y_{00101}=\frac{1}{7}+Y_{11010},\\
Y_{00110}&=\frac{3}{7}-Y_{11010}-Y_{11100},\qquad
Y_{11001}=\frac{2}{7}-Y_{11010}-Y_{11100},
\end{aligned}
\]
with $Y_{11010},Y_{11100}\ge 0$ and $Y_{11010}+Y_{11100}\le\tfrac{2}{7}$.

\paragraph{Two free parameters, but the same marginal vector.}
Let
\(
 u:=Y_{11010},\ v:=Y_{11100}
\).
A direct computation using \eqref{eq:Zi-diagonal} shows that for \emph{all} $(u,v)$ in the triangle
\(
 u\ge 0,\ v\ge 0,\ u+v\le \tfrac{2}{7}
\), the shared marginals are constant:
\[
\bigl(\langle Z_1\rangle,\dots,\langle Z_5\rangle\bigr)=\left(\frac{3}{7},\frac{3}{7},-\frac{1}{7},-\frac{1}{7},-\frac{1}{7}\right).
\]
Consequently
\[
\Sigma_Z=\sum_{i=1}^5\langle Z_i\rangle^2
=2\left(\frac{3}{7}\right)^2+3\left(\frac{1}{7}\right)^2
=\frac{3}{7}
\]
is also fixed.

\paragraph{Interpretation.}
Although the feasible set in $(X,Y)$ has dimension $2$ (the triangle in $(u,v)$), these parameters only change \emph{how} the common marginal vector is expressed as a convex combination of vertices in $V_1$; the vector itself is fixed. Equivalently, $\mathrm{conv}(V_0)\cap\mathrm{conv}(V_1)$ is a single point, and the $(u,v)$ freedom is the ``fiber'' of decompositions realizing that point.

\subsection{Practical test for ``varying marginals'' families}

Given supports $C_0,C_1$, one clean way to determine whether the shared marginals can vary is:
\begin{itemize}
\item Solve the feasibility LP (equalities + nonnegativity). If infeasible, stop.
\item Optimize a coordinate of the marginal vector, e.g. maximize and minimize $\langle Z_1\rangle$ subject to the same constraints. Because $\langle Z_1\rangle$ is linear in $(X,Y)$, this is an LP.
\item If the maximum differs from the minimum, then $\mathrm{conv}(V_0)\cap\mathrm{conv}(V_1)$ has positive dimension and you can build an explicit one-parameter family by convex-combining the two extremal solutions.
\end{itemize}
This criterion matches the empirical distinction between the worked $m=7$ example (fixed marginals) and Family~A (varying marginals).


\end{document}
